\begin{abstract}
This teaching outline presents a research project on the construction, visibility, and state mobilization of lethal violence indicators in Brazil, covering reference years from 2016 to 2024. The research asks how choices in constructing and presenting these indicators relate to their institutional and media visibility and their mobilization in federal processes of attention and prioritization. Its objective is to analyze these relationships through a documentary, longitudinal, mixed-methods design. The theoretical foundation connects commensuration, comparisons across arenas of attention, and the political mobilization of indicators. The proposed approach distinguishes administrative records, statistical consolidation and analytical products, institutional communication, journalistic selection, and federal legislative and executive instruments. It preserves an original data layer and a separate analytical layer and does not automatically treat differences between sources as errors. News coverage will be examined in G1, Folha de S.Paulo, and O Estado de S.Paulo, without claiming statistical representativeness of the entire press. The expected technical output is a protocol for indicator traceability, comparability, and responsible communication. This document demonstrates the IDP LaTeX template, citations, and cross-references. It does not report empirical findings or infer causality from documentary correspondence.
\end{abstract}
